%!TeX program = pdflatex
\documentclass[acmtog,review,anonymous]{acmart}

\AtBeginDocument{\providecommand\BibTeX{{Bib\TeX}}}
\citestyle{acmauthoryear}
\setcopyright{none}
\settopmatter{printacmref=true}
\renewcommand\footnotetextcopyrightpermission[1]{}
\acmSubmissionID{TBD}
\acmJournal{TOG}
\acmVolume{XX}
\acmNumber{X}
\acmArticle{XX}
\acmYear{2026}
\acmMonth{9}

\usepackage{amsmath}
\usepackage{amssymb}
\usepackage{booktabs}
\usepackage{graphicx}
\usepackage{xcolor}
\usepackage{xspace}

\graphicspath{{submissions/2026-09/figures/}{figures/}{../figures/}}

\newcommand{\method}{Socratic Method\xspace}
\newcommand{\model}{M7-XR\xspace}
\newcommand{\todo}[1]{\textcolor{red}{\textbf{TODO: #1}}}
\newcommand{\placeholder}[3]{%
  \fbox{\begin{minipage}[c][#2][c]{#1}
  \centering\small\textbf{FIGURE RESERVED}\\[0.6em]
  \texttt{\detokenize{#3}.pdf} or \texttt{\detokenize{#3}.png}
  \end{minipage}}}
\newcommand{\figureasset}[3]{%
  \IfFileExists{submissions/2026-09/figures/#1.pdf}%
    {\includegraphics[width=#2]{submissions/2026-09/figures/#1.pdf}}%
    {\IfFileExists{submissions/2026-09/figures/#1.png}%
      {\includegraphics[width=#2]{submissions/2026-09/figures/#1.png}}%
      {\placeholder{#2}{#3}{#1}}}}

% Generated from the live v31 history by scripts/snapshot_live_v31.py.
\newcommand{\observedSamples}{8192}
\newcommand{\observedSamplesText}{8,192}
\newcommand{\bestSamples}{7168}
\newcommand{\bestSamplesText}{7,168}
\newcommand{\bestMacro}{0.58293}
\newcommand{\releaseSamples}{8192}
\newcommand{\releaseSamplesText}{8,192}
\newcommand{\releaseThreshold}{0.45}
\newcommand{\durationRows}{%
1,024 & .57944 & .57136 & +.01968 & +.01917 \\
2,048 & .57886 & .57086 & +.01910 & +.01868 \\
3,072 & .58037 & .57253 & +.02062 & +.02037 \\
4,096 & .58022 & .57196 & +.02047 & +.02024 \\
5,120 & .57873 & .57014 & +.01897 & +.01813 \\
6,144 & .58104 & .57329 & +.02128 & +.02124 \\
7,168 & \textbf{.58293} & \textbf{.57537} & \textbf{+.02317} & \textbf{+.02276} \\
8,192 & .58037 & .57226 & +.02061 & +.02031 \\
}


\begin{document}

\title{The Socratic Method: Cross-Resolution Supervision for Papyrus Surface Segmentation}

\begin{abstract}
Coarse-resolution papyrus surface segmentation fails in two opposing ways:
faint sheets disappear, while nearby wraps merge into blobs.
We present a teacher-student method for Herculaneum papyrus data in which an
approximately $2.399\,\mu$m native-fine model interrogates the released
$9.362\,\mu$m M7 segmenter.
The resulting student remains an ordinary, single M7 residual-encoder nnU-Net
at inference.
During training, projected soft occupancy is combined with medial-crest recall,
teacher-background separation, M7 function preservation, a global parameter
trust region, and dynamic widest-path connectivity.
A separate additive 2D curve fitter reconnects short, geometrically credible
gaps after segmentation.
Our completed duration study improves both held-out scrolls over released M7 at
every measured interval through \observedSamplesText\ samples, with a best
calibrated macro-scroll Dice of $\bestMacro$ at \bestSamplesText\ samples.
An independent registered-morphology and blind anti-blob review selects the raw
\releaseSamplesText-sample student at threshold $\releaseThreshold$. It passes
the six-cube PHerc1447 anti-blob gate and matches teacher component count on 15
of 16 locked PHerc0139 slices; the remaining scalar exception and the difference
between perceptual and overlap-based rankings are reported explicitly.
\end{abstract}

\begin{CCSXML}
<ccs2012>
 <concept>
  <concept_id>10010147.10010257.10010293.10010294</concept_id>
  <concept_desc>Computing methodologies~Neural networks</concept_desc>
  <concept_significance>500</concept_significance>
 </concept>
</ccs2012>
\end{CCSXML}
\ccsdesc[500]{Computing methodologies~Neural networks}
\keywords{volumetric segmentation, knowledge distillation, papyrus,
topology-aware learning, curve completion}

\begin{teaserfigure}
  \centering
  \figureasset{teaser}{0.98\linewidth}{3.0in}
  \caption{One raw checkpoint answers two opposing failure modes. Both rows
  begin with the identical CT crop and released M7 state-of-the-art baseline.
  The top row shows blind PHerc1447 de-blob behavior; the bottom row adds the
  projected fine teacher and shows longitudinal growth on locked PHerc0139.
  Every M7-XR panel uses the selected \releaseSamplesText-sample student at
  $t=\releaseThreshold$, without blending.}
  \Description{Two-row comparison. Each row shows CT, released M7, and the raw
  M7-XR student; the locked longitudinal-growth row also shows its projected
  fine-teacher target.}
  \label{fig:teaser}
\end{teaserfigure}

\maketitle

\section{Introduction}
\label{sec:introduction}

\begin{quote}
\centering\small
``I neither know nor think that I know.''\\[-0.2em]
---Socrates, in Plato's \emph{Apology} 21d~\cite{platoApology21d}

\textsc{Socrates:} I am wiser than this man; he fancies he knows something,
although he knows nothing---\\
\textsc{Darryl, Socrates' friend:} \emph{fuck him up socrates}\\[-0.2em]
---leon (@leyawn)~\cite{leyawn2015socrates}
\end{quote}

The objects behind this work are the Herculaneum papyri, a library of scrolls
carbonized by the AD~79 eruption of Mount Vesuvius.
Carbonization preserved the rolls but left them too brittle to open and the
writing trapped inside tightly wound coils.
Synchrotron micro-CT can record their internal structure without touching the
object; recovering a page then requires locating the papyrus sheet in three
dimensions and \emph{virtually unwrapping} it to a plane
\cite{seales2016damage,vesuviusChallenge}.
A scroll is one continuous sheet, metres long, visiting many adjacent turns.
Consequently, two pieces separated by only a few CT voxels may be distant along
the page, while a short gap in the volume may interrupt one continuous line of
text.

Segmenting the papyrus surface is the first geometric commitment in a
scroll-unwrapping pipeline.
A missed surface interrupts a sheet that may otherwise be traceable over a long
distance.
An overgrown prediction is at least as dangerous: it can weld nearby turns into
one component and make a locally plausible surface globally wrong.
A useful segmentation method must therefore recover faint structure and remove
false connections at the same time.
An overlap score alone does not adequately describe this tension.

The Vesuvius data provides a second, complementary source of information.
A model operating on finer CT can resolve local evidence that is ambiguous after
downsampling, while the deployed segmentation and geometry stack still consumes
the coarser grid.
We treat this mismatch as a cross-resolution dialogue.
The fine model proposes soft occupancy and medial evidence; the released M7
model contributes a function that is already effective; explicit loss terms
expose where those views are inconsistent.
The student must reconcile them without moving far from M7.

This is the origin of our title.
In the Socratic account, a speaker who professes little knowledge reveals
inconsistencies in an expert's position and thereby points toward a better one.
Our fine teacher is similarly not deployed as a superior replacement.
It asks local questions of the expert coarse model.
The answer is a revised M7 that stands alone at inference.

We present four contributions:
\begin{itemize}
  \item a hash-pinned $2\,\mu$m-to-$9\,\mu$m teacher-supervision recipe that
  can be executed independently of our geometry repository;
  \item an objective that balances soft occupancy, de-blob separation, medial
  recall, dynamic connectivity, and two forms of M7 preservation;
  \item a duration-ladder protocol that retains each declared observation rather
  than assuming the longest run is best; and
  \item an independently measured, additive 2D curve fitter that reconnects
  only short gaps able to answer geometric and cross-plane objections.
\end{itemize}

The 8,192-sample duration experiment and the independent morphology review are
complete. We report every committed interval, including the fact that the
best-Dice interval is not the released interval. The selected artifact is the
raw 8,192-sample student at threshold 0.45; its known component-count exception,
undergrowth cases, and gate-proxy failures remain visible in the record.

\section{Related Work}
\label{sec:related}

\paragraph{Segmentation.}
U-Net architectures and their self-configuring nnU-Net variants are widely used for dense volumetric segmentation
\cite{ronneberger2015unet,isensee2021nnunet}. Our deployed network follows the released M7 residual-encoder nnU-Net from the
Vesuvius Challenge project exactly; we alter its supervision and optimization rather than adding an inference-time branch.
This preserves compatibility with the existing coarse-volume pipeline.

\paragraph{Teacher supervision.}
Knowledge distillation transfers information from a teacher into a student by matching soft outputs or intermediate functions \cite{hinton2015distilling}.
Most such methods assume teacher and student predictions share a sampling grid. Our teacher observes a finer CT volume, so its occupancy and validity must be
pulled back into the student's physical frame. The resulting targets represent partial volume rather than only a hard class.

\paragraph{Topology-aware Objectives.}
Skeleton and centerline losses preserve thin connected structures that can be under-represented by region overlap metrics
\cite{shit2021cldice}. The medial axis transform (MAT) instead describes a
shape by centers and the radii of their maximally inscribed disks or spheres.
LSMAT casts this construction as a least-squares problem to make the resulting
representation less brittle under noisy or perturbed boundaries
\cite{rebain2019lsmat}. This center--radius factorization is especially useful
for papyrus: motion along the medial set extends a sheet, whereas increasing
radius thickens it toward a neighboring wrap. We use that distinction to derive
thin training corridors and M7 contact pins. Unlike a conventional skeleton
overlap loss, our final connectivity term does not prescribe every voxel of one
centerline. It raises the bottleneck of whichever admissible medial path the
student currently supports. The atlas is constructed only from fully owned
training events and never from held-out data.

\paragraph{Curve completion.}
Classical thinning extracts one-pixel skeletons from binary images
\cite{zhang1984fast}. Curve completion can then be posed as endpoint pairing under geometric
continuity. For wound papyrus, however, a locally smooth completion can cross from one wrap
to the next. Our postprocess therefore combines 2D tangents with an umbilicus-aware radial
test, neighboring-plane corridors, explicit anti-merger and crossing gates, and
connection persistence across slices.

\paragraph{Vesuvius models and data.}
Our experiments build on the public Vesuvius Challenge CT data, the released M7
surface model, and a Villa native-fine teacher
\cite{vesuviusChallenge,m7Artifact,villaTeacher}.
The final paper will replace these artifact-level entries with the authoritative
model and dataset citations associated with the public release.

\section{Overview}
\label{sec:overview}

Fig.~\ref{fig:overview} summarizes our pipeline.
We first run a native-fine teacher in its approximately $2.399\,\mu$m CT frame.
Its probabilities, support, and medial crest are projected into the
$9.362\,\mu$m grid as soft training fields (Sec.~\ref{sec:crossres}).
We initialize every candidate from the exact released M7 checkpoint and train
all parameters under a small global trust region. The output is one ordinary M7 network.
At inference we use only its raw class-1 probability; there is no teacher and no
M7 probability blend.

At the selected threshold, an optional C postprocess operates on each axial
slice (Sec.~\ref{sec:fitter}). It proposes Bezier completions between skeleton endpoints, but paints a join
only when local curve geometry, radial displacement, neighboring slices, and
connection-track support agree. Segmentation training and curve completion are measured separately so an
improvement in one stage cannot hide a failure in the other.
All segmentation qualification numbers in Sec.~\ref{sec:results} use the raw
student without morphology or line fitting.

\begin{figure*}[t]
  \centering
  \figureasset{figure_1_method_overview}{0.96\linewidth}{2.05in}
  \caption{Method overview. The final figure will separate training-only
  components (fine teacher, projection, auxiliary losses, and frozen M7 anchor)
  from deployed components (raw student and optional 2D fitter). It will also
  show that the parameter trust-region projection acts after every optimizer
  update.}
  \Description{Reserved schematic of training-only and deployed pipeline components.}
  \label{fig:overview}
\end{figure*}

Two engineering details prevent accidental leakage or model substitution. First, the dynamic connectivity atlas is built only from boxes in the training
manifest; held-out morphology is used only for evaluation. Second, the M7 initializer is accepted only if both its byte size and SHA-256
match the released artifact. A new candidate may resume its own interrupted optimizer state, but may not use
an earlier student as learned initialization.

\section{Cross-Resolution Supervision}
\label{sec:crossres}

\subsection{Student Architecture and Data}

Our student is a six-stage 3D residual-encoder nnU-Net with feature counts
$(32,64,128,256,320,320)$, one CT input channel, and two output classes.
It has $102{,}349{,}770$ trainable parameters.
The frozen training schedule contains 4,096 deterministic $192^3$ patches from
PHerc0139 and a single native-fine teacher record. We traverse this schedule
twice for 8,192 samples.
The small corpus reflects the time available before this progress submission,
not the extent of the registered data. Our broader pair plan contains three
additional non-prize training scrolls---PHerc0332, PHerc1667, and PHercParis4---
with 62,500 candidate origins allocated to each. None of the four training
scrolls is eligible for a Herculaneum prize.

Our held-out validation manifest contains registered PHerc0814 and PHerc1451
patches. We evaluate every 1,024 samples and retain model-only checkpoints at 1,024,
2,048, 4,096, and 8,192 samples. Because the measured duration response is not monotonic
(Sec.~\ref{sec:results}), we do not assume the final checkpoint is best.

\subsection{Fine-Teacher Occupancy and Medial Geometry}

The fine teacher provides two complementary targets in the student's physical
frame. The first is soft occupancy $q\in[0,1]$, obtained by an anti-aliased
pullback from the $2.399\,\mu$m grid. It estimates how much of a $9.362\,\mu$m
voxel is occupied by papyrus. The second is a binary medial crest $c$, projected
with an OR/max operator so that the existence of a fine-scale center survives
downsampling. Occupancy and crest have independent validity masks.

This separation follows a concrete failure of occupancy-only supervision. A
thin but real continuation commonly occupies only 25--45\% of a coarse voxel.
Soft cross-entropy is correctly calibrated when it drives such a voxel toward
$p=0.25$--$0.45$, but the result may remain below the segmentation threshold.
Reducing background separation moved both the center and its shoulders upward,
recovering length while reintroducing thick blobs. We therefore retain $q$ as
an amount-of-material target and use $c$ as an independent sheet-existence
target.

\paragraph{LSMAT-style center--radius representation.}
LSMAT represents a shape using medial centers and their maximally inscribed
radii rather than boundary samples alone~\cite{rebain2019lsmat}. We use that
center--radius view while retaining the released Villa teacher's center
construction. For each native-fine axial slice $T_z$, the teacher skeletonizes
the slice in 2D; the stacked centerlines receive one 3D morphological closing
and are intersected with the original foreground. We pair every surviving
center $u\in C_z$ with a physical 2D Euclidean-distance radius
\begin{equation}
 r_z(u)=\min_{y\notin T_z}\lVert u-y\rVert_{\mathrm{phys}},
 \qquad
 \widehat T_z=\bigcup_{u\in C_z} B^{\circ}(u,r_z(u)).
 \label{eq:center-radius}
\end{equation}
The disks are open because the stored distance reaches the nearest background
voxel center. Slice-wise radii are deliberate: pairing 2D centers with a 3D
distance transform would mix disks and spheres and collapse the desired medial
surface toward a curve. We audit reconstruction, component count, false-positive
girth, and halo stability before using these centers for supervision.

The representation gives us two independent geometric coordinates. Moving
along $C_z$ extends the papyrus sheet tangentially; increasing $r_z$ grows it
radially toward another wrap. That distinction is the reason for introducing
the LSMAT view: the connectivity loss should act along the first coordinate,
while the separation loss and radius screens police the second.

\paragraph{From medial centers to pins.}
We construct events only on training slices with at least 95\% valid teacher
support. Components smaller than 12 voxels are removed. A slice qualifies when
one thin teacher component $H_e$ contacts at least two disconnected M7
components $M_{e,i}$ and contains at least two teacher voxels absent from the
one-voxel M7 neighborhood. Both teacher and M7 must have maximum in-plane EDT
radius at most three voxels, and no more than 2\% of their foreground may
survive a two-voxel erosion; these checks reject blobs before they can define a
positive connectivity target.

Let $C_e$ be the projected medial centers inside $H_e$, and let $D_r$ denote
radius-$r$ dilation. We choose the smallest $d\in\{0,1,2,3\}$ for which
\begin{equation}
 \Omega_e = H_e\cap D_d(C_e),
 \qquad
 A_{e,i}=\Omega_e\cap D_2(M_{e,i})
 \label{eq:pins}
\end{equation}
contains every nonempty contact set $A_{e,i}$ in one connected component. The
sets $A_{e,i}$ are the \emph{pins}: they are where distinct released-M7
fragments meet the LSMAT-style teacher-medial corridor, not arbitrary endpoints
chosen from a student prediction. Corridor voxels within one voxel of M7, plus
all pin voxels, form the free-anchor set $F_e$. Events touching their owning
tile boundary are discarded, so every retained corridor is transformed and
loaded as one unit.

\subsection{Objective}

Let $p$ be the student's surface probability, $q$ the projected fine-teacher
occupancy, $c$ its binary medial-crest field,
$C=\{x:c_x=1\}$ the corresponding valid crest set, and $p_0$ the released M7
probability.
Our loss is
\begin{align}
\mathcal{L} ={}& \mathcal{L}_{\mathrm{CE}}(p,q)
 + 0.25\mathcal{L}_{\mathrm{Dice}}(p,q)
 + \mathcal{L}_{\mathrm{crest}}(p,c) \nonumber\\
&+2\mathcal{L}_{\mathrm{sep}}(p,q)
 +0.5\mathcal{L}_{\mathrm{KL}}(p,p_0)
 +\mathcal{L}_{\mathrm{pres}}(p,p_0) \nonumber\\
&+0.03125\mathcal{L}_{\mathrm{conn}}(p).
\label{eq:objective}
\end{align}
All terms are evaluated only where their corresponding validity fields permit.

\paragraph{Soft Occupancy and Crest Recall.}
Our cross-entropy and Dice targets retain fractional occupancy created by
fine-to-coarse projection. For a nonempty valid crest $C$, the additional term is
\begin{equation}
 \mathcal{L}_{\mathrm{crest}}
 = 1-\frac{\sum_{x\in C}p_x+1}{|C|+1}.
 \label{eq:crest-loss}
\end{equation}
We average Eq.~\ref{eq:crest-loss} over crest-bearing samples rather than over
all voxels. The normalization prevents sparse medial evidence from disappearing
into the background count, and the additive one is the smoothing constant used
by the released teacher loss. Because only crest voxels receive its gradient,
the term asks for longitudinal sheet evidence without defining a thick target.

\paragraph{De-Blob Separation.}
Thick predictions need explicit shrink pressure. An M7 blob can overlap the
fine sheet and still score well against an anti-aliased occupancy target.
Let $V$ and $V_c$ denote occupancy and crest validity. We form the thin positive
seed
\begin{equation}
 P=C\cup\bigl(\{q\geq0.5\}\cap(V\setminus V_c)\bigr),
 \label{eq:separation-positive}
\end{equation}
using the crest wherever it is known and falling back voxelwise to hard
occupancy only where it is not. We dilate $P$ by two voxels and retain only
occupancy-valid locations where the teacher is confidently background:
\begin{equation}
  S = \bigl(\operatorname{dilate}_2(P)\setminus P\bigr)
      \cap V\cap \{q\leq0.1\}.
  \label{eq:separation-shell}
\end{equation}
The separation loss is the mean background cross-entropy
\begin{equation}
  \mathcal{L}_{\mathrm{sep}}
  = -\frac{1}{|S|}\sum_{x\in S}\log(1-p_x).
  \label{eq:separation-loss}
\end{equation}
We weight this term by 2 in Eq.~\ref{eq:objective}. Because $S$ is a local fence around supported medial
structure rather than arbitrary background, the term attacks thickness and
short cross-wrap welds while the soft occupancy and crest terms remain free to
grow a faint sheet elsewhere.

\paragraph{M7 function and positive preservation.}
The teacher is locally more informative, not globally complete. We therefore
freeze a copy of released M7 during training and use its probability $p_0$ in
two differently masked terms. Let
$H=C\cup(\{q\geq0.5\}\cap V)$, $h_x=\mathbf{1}[x\in H]$, and
$m_x=\mathbf{1}[p_{0,x}\geq0.5]$. The KL mask is
\begin{align}
 K={}&\bigl((\overline V)\setminus D_2(H)\bigr) \nonumber\\
 &{}\cup\{x\in V:h_x=m_x\ \wedge\ (q_x\leq0.1\ \vee\ h_x=1)\}.
 \label{eq:kl-mask}
\end{align}
Thus unknown space is anchored except for a radius-two corridor around known
teacher positives. In known space, M7 is anchored only where its hard decision
agrees with a confident teacher decision; ambiguous partial occupancy is not
misclassified as agreement on background. We use
\begin{equation}
 \mathcal{L}_{\mathrm{KL}}
 =\frac{1}{|K|}\sum_{x\in K}
 D_{\mathrm{KL}}\!\left(\operatorname{Ber}(p_{0,x})\,
 \middle\|\,\operatorname{Ber}(p_x)\right).
 \label{eq:m7-kl}
\end{equation}

KL preservation alone can still allow a correct M7 foreground fragment to
shrink under the separation term. We therefore define
$R=V\cap D_2(H)\cap\{p_0\geq0.5\}$ and add the positive-only loss
\begin{equation}
 \mathcal{L}_{\mathrm{pres}}
 =-\frac{1}{|R|}\sum_{x\in R}\log p_x.
 \label{eq:m7-preservation}
\end{equation}
The two terms have separate roles: Eq.~\ref{eq:m7-kl} retains M7's foreground
and background function away from justified edits, while
Eq.~\ref{eq:m7-preservation} prevents de-blobbing from deleting M7-supported
sheet near the teacher. Teacher disagreement remains free to correct M7.

\paragraph{Dynamic medial connectivity.}
Separation supplies the desired shrink pressure, but without a complementary
topological term it can also sever the same faint longitudinal connections we
want to retain. Our first LSMAT axial loss connected the same pins with a
minimum-off-axis spanning tree, then raised the weakest 10\% of voxels on that
fixed route toward probability 0.2. This proved that a thin, center-biased route
could be supervised without using held-out geometry, but it also froze one of
many equally valid paths. The lowest weight only tied the incumbent; larger
weights monotonically reduced ordinary validation quality, and the strongest
weight also reduced anti-blob passes. The final loss therefore retains the
audited LSMAT corridor and pins but makes the path existential.

For event $e$, define the path capacity
\begin{equation}
 a_e(x)=
 \begin{cases}
 1, & x\in F_e,\\
 p_x, & x\in\Omega_e\setminus F_e,\\
 0, & x\notin\Omega_e.
 \end{cases}
 \label{eq:path-capacity}
\end{equation}
Taking $A_{e,0}$ as a source pin, initialize
$r_e^0(x)=\mathbf{1}[x\in A_{e,0}]$ and iterate
\begin{align}
 r_e^{t+1}(x)=\mathbf{1}[x\in\Omega_e]\max\!\Bigl(&r_e^t(x),
 \min\bigl(a_e(x),\max_{y\in\mathcal{N}_{26}(x)}r_e^t(y)\bigr)\Bigr).
 \label{eq:widest-recurrence}
\end{align}
This max--min recurrence computes the greatest bottleneck capacity among paths
of at most $t$ steps. After 96 steps, the event score and loss are
\begin{align}
 b_e & = \min_{i>0}\max_{x\in A_{e,i}} r_e^{96}(x), \nonumber\\
 \mathcal{L}_{\mathrm{conn}}
 & = \frac{1}{|E|}\sum_{e\in E}
 \max\bigl(0,\operatorname{logit}(0.2)-\operatorname{logit}(b_e)\bigr).
 \label{eq:connectivity-loss}
\end{align}
Each event receives one vote regardless of corridor size or pin count. Max/min
back-propagation places pressure on the current best path's limiting capacity,
not on every corridor voxel. Free M7 anchors have unit capacity and no direct
gradient; no connectivity gradient exists outside $\Omega_e$. The loss becomes
exactly zero when every pin is reachable through some path with bottleneck at
least 0.2.

The frozen atlas contains 149 fully owned events. Ninety-nine rows in the exact
schedule encounter an event, 477 rows are eligible, and the longest audited
event needs 44 of the allocated 96 propagation steps. Of these events, 135 have
two pins, 13 have three, and one has four; 122 need no crest dilation and 27 need
one voxel. The loss is evaluated only at the full-resolution output and has
weight 0.03125. Together, the separation shell and existential path constraint
express the two sides of de-blobbing: remove unsupported girth without deleting
sheet existence.

\paragraph{Deep supervision.}
The occupancy, crest, separation, KL, and preservation terms are applied at the
standard nnU-Net decoder outputs with normalized weights proportional to
$1,1/2,1/4,\ldots$ and zero weight on the coarsest output. Soft occupancy is
resized trilinearly. The sparse crest is max-pooled, while its validity is true
only if every contributing fine cell is valid. Connectivity is intentionally
excluded from lower-resolution outputs so resampling cannot broaden its medial
corridor.

\begin{figure*}[t]
  \centering
  \figureasset{figure_crossres_supervision}{0.96\linewidth}{2.0in}
  \caption{Actual PHerc0139 training data at dynamic event 138. Projected soft
  occupancy, its sparse medial crest, and the radius-two separation shell encode
  center/radius geometry without rewarding girth. The lower panels show the M7
  fragments, teacher-medial corridor, component-contact pins, and free anchors
  supplied to the max--min connectivity recurrence.}
  \Description{Training slice with projected teacher occupancy, medial crest,
  de-blob shell, M7 fragments, and dynamic-connectivity metadata.}
  \label{fig:supervision}
\end{figure*}

\subsection{Optimization and Trust Region}

We use Adam with constant learning rate $2\times10^{-5}$, betas $(0.9,0.999)$,
$\epsilon=10^{-8}$, and zero weight decay.
The batch size is three with eight-step gradient accumulation and bfloat16
autocast.
Seed 1203 fixes the sample order and augmentation sequence.

After each optimizer update, we project the complete parameter vector $\theta$
into a relative-$L^2$ ball around released parameters $\theta_0$:
\begin{equation}
\theta \leftarrow \theta_0 +
\min\!\left(1,\frac{r\lVert\theta_0\rVert_2}
{\lVert\theta-\theta_0\rVert_2}\right)(\theta-\theta_0),
\qquad r=0.0027535422.
\label{eq:trust}
\end{equation}
The projection is active on essentially every observed update.
It is thus part of the effective method rather than an inactive guardrail.

\section{Additive 2D Curve Fitting}
\label{sec:fitter}
Following re-training, we introduce a downstream postprocess that repairs short omissions
in the thresholded predictions. It is additive only: accepted joins set foreground pixels, but never erase
the model output. This restriction makes every intervention easy to localize and prevents a
gap-repair stage from quietly becoming a second segmentation method.

For each $z$-slice, we apply Zhang--Suen thinning~\cite{zhang1984fast}, remove short spurs, label
the 8-connected skeleton and mask components, and extract curve endpoints with
outward tangents. We then snap endpoints to the mask edge so distances describe the visible gap rather
than skeleton retraction. Endpoints clipped by the volume or missing data are retained for audit but may
not be matched.

We score endpoint pairs using distance, facing, adjacent-plane evidence, and cross-slice track support.
Before a pair enters deterministic greedy matching it must satisfy the following requirements:

\begin{enumerate}
  \item it does not close the same skeleton or mask component;
  \item its tangents face the gap and are sufficiently opposed;
  \item its reach is at most 21 pixels, while the calibrated safe regime at six
  pixels or less requires less supporting evidence;
  \item if an umbilicus is available, its radial displacement is at most three
  pixels;
  \item its cubic Bezier has a bounded arc-to-chord ratio and does not approach
  a third component;
  \item the path corridor does not cross another sheet on nearby slices; and
  \item it neither crosses nor touches an already accepted join, and neither
  endpoint has taken a better match.
\end{enumerate}

A first matching pass seeds endpoint and connection tracks.
A second pass follows each new track through $z$, uses connection confidence, and retains
only joins supported across at least three slices.
We then paint a cubic Bezier disk stroke and record every added pixel.

\begin{figure*}[t]
  \centering
  \figureasset{figure_3_line_fitter_examples}{0.96\linewidth}{2.9in}
  \caption{Measured line-fitter behavior on the released M7 baseline. Left:
  representative M7 masks before and after additive accepted joins. Center:
  the same connection persists across neighboring $z$-slices. Right:
  registered fine-CT evidence distinguishes a supported connection from a
  required separation. The calibration accepted 481 of 772 proposals, added
  12,697 pixels, reached 92.5\% registered fine-CT precision, and produced no
  observed full-turn fusion. Cyan is released M7 and green is fitted addition.}
  \Description{Released M7 before-and-after curve fitting, cross-slice track
  persistence, connected and separate fine-CT validation examples, and measured
  calibration statistics.}
  \label{fig:fitter}
\end{figure*}

\section{Results}
\label{sec:results}

\subsection{Evaluation Protocol}

We calibrate thresholds from 0.25 through 0.60 in increments of 0.01 on
registered PHerc0814 and PHerc1451 patches.
Checkpoint ranking uses calibrated macro-scroll Dice with a minimum per-scroll
gain guard against released M7.
Final qualification evaluates the raw student at five durations
($1{,}024$, $2{,}048$, $3{,}072$, $4{,}096$, and $8{,}192$ samples) and
thresholds 0.38--0.50, with 0.25 retained as a failure control. Its evidence is
16 locked PHerc0139 slices and six blinded PHerc1447 cubes. These locations are
evaluation-only and spatially separated from the PHerc0139 training boxes.
Promotion requires anti-blob behavior, two-sided morphology and coverage, and
human review; no morphology, M7 blend, or line fitter is included in these
model measurements.
An earlier raw v29 candidate failed the PHerc1447 gate; we do not consider it a
viable model.

\subsection{Duration Study}

Table~\ref{tab:duration} reports every committed validation interval through the
completed \observedSamplesText-sample schedule.
All intervals improve both held-out scrolls over released M7.
The response is shallow and non-monotonic; the best is the
\bestSamplesText-sample interval at $\bestMacro$ calibrated macro-scroll Dice.
This behavior justifies retaining the full ladder instead of exposing only the
longest checkpoint.

\begin{table}[t]
\caption{Completed v31 duration ladder. ``Cal.'' is calibrated macro-scroll Dice;
``$t=.40$'' uses a fixed threshold; ``gain'' is calibrated macro gain over the
M7 initial model. All calibrated thresholds equal the 0.25 search boundary.}
\label{tab:duration}
\centering
\small
\begin{tabular}{rrrrr}
\toprule
Samples & Cal. & $t=.40$ & Gain & Min. gain \\
\midrule
\durationRows
\bottomrule
\end{tabular}
\end{table}

The trust-region projection is active for 97.7\% of optimizer updates in the
first interval and 100\% thereafter.
The method therefore reaches its best observed validation score while remaining
on the boundary of its allowed parameter neighborhood.

Every interval selects threshold 0.25, the lower edge of our sweep.
The calibrated scores therefore compare duration, but do not determine the
operating point: the apparent optimum is censored and 0.25 visibly refills the
blobs that the independent gate is designed to catch.

\subsection{Morphology-Gated Selection}

Joint review selects the raw \selectedSamplesText-sample checkpoint at
$T=\selectedThreshold$. At $T=0.42$ it obtains 16/16 exact component-count
matches and already passes the PHerc1447 anti-blob gate. We nevertheless select
0.45 for its cleaner useful-structure/de-blob balance. At that operating point,
15/16 locked slices have exact component counts; the remaining case is a known
scalar ambiguity in which the intended structures are present but one feature
splits while two strokes touch.

On the six PHerc1447 cubes, the selected model has foreground ratio 0.950508
relative to the v15 anti-blob reference, reference recall 0.911324, zero
interior-regression cubes, zero thickness-regression cubes, and a hard
anti-blob pass. Human review finds useful line recovery in difficult views
without restoration of the released-M7 blobs, although several undergrown
views remain explicit failure watches. The full checkpoint hash is in the
machine-readable release recipe.

\paragraph{Independent perceptual audit.}
We also run FLIP~\cite{andersson2020flip} at 67.0206 pixels per degree on the
exact 1,120 displayed candidate/teacher pairs. At $T=0.45$, the 8,192-sample
model improves FLIP on 12/16 locked slices relative to 2,048 samples; total
teacher-only erosion falls from 4,283 to 4,270 pixels and false additions from
1,201 to 1,155. The audit is not unanimous: mean FLIP narrowly prefers 4,096
samples (0.232073 versus 0.233310), while weighted-median pooling narrowly
prefers 8,192 (0.599964 versus 0.600288). We retain this disagreement. FLIP
rewards teacher resemblance, whereas the blind gate and human review also
penalize the teacher-like girth that can weld adjacent wraps.

\begin{figure*}[t]
  \centering
  \figureasset{figure_2_registered_examples}{0.96\linewidth}{2.2in}
  \caption{Reserved registered comparison. Columns will show CT, released M7,
  projected fine-teacher supervision, the selected 8,192-sample output at
  $T=0.45$, and two-sided
  add/remove errors. Rows will cover PHerc0814, PHerc1451, and the blinded
  PHerc1447 gate, with identical crops and thresholds clearly reported.}
  \Description{Reserved registered comparison across held-out and blind scrolls.}
  \label{fig:registered}
\end{figure*}

\subsection{Line-Fitter Evaluation}

In a calibrated $4\!\times\!5\!\times\!5$ study, the line fitter accepted 481
joins and added 12,697 pixels.
Registered high-resolution precision was 92.5\% overall and 98.4\% for joins no
longer than six pixels.
We observed no full-turn fusion and measured a 21\% reduction in atlas overlap.
A larger 21-cube census measured 98.3\% weighted precision.

Adjacent high-resolution connectivity is useful confirmation but is not, by
itself, proof that a join remains on one wrap.
The umbilicus-aware radial gate is our primary cross-wrap safeguard.
The postprocessor evaluation therefore reports false joins by gate, distance, and
radial regime rather than only an aggregate precision.

\subsection{Interpretation}

The duration experiment establishes consistent positive held-out gains under a
saturated M7 trust constraint, while the completed morphology audit establishes
a release candidate. Their different choices are the central result:
overlap-only calibration prefers 7,168 samples and threshold 0.25, but the
de-blob-aware evidence selects 8,192 and 0.45. Table~\ref{tab:duration} is
generated only from committed run artifacts, and the selection record preserves
the independent locked, blind, and perceptual evidence rather than collapsing
them into one score.

\section{Limitations}
\label{sec:limitations}

Our fitting schedule is PHerc0139-only and may encode scroll-specific material
or acquisition properties.
Fine-teacher predictions are model-derived soft supervision rather than human
ground truth.
Two held-out calibration scrolls, 16 locked slices, and six blind cubes provide
guards but do not establish broad robustness.
The M7 trust projection is saturated, and its radius was selected under an
earlier objective rather than the final loss in Eq.~\ref{eq:objective}.
Our overlap-calibration sweep also fails to bracket its optimum; threshold 0.45
is a morphology-gated operating choice, not the maximizer of that curve. Human
review is necessarily less reducible than a scalar and should be repeated by
independent readers before a public model claim.

The line fitter has a different failure mode.
A completion may be locally smooth and persistent across slices but still join
the wrong wrap near the core.
Radial, third-component, and cross-sheet gates reduce this risk but cannot prove
global correctness.
The postprocess must remain independently measured and its additions must remain
recoverable during qualification.

\begin{figure}[t]
  \centering
  \figureasset{figure_4_failure_cases}{0.94\columnwidth}{1.45in}
  \caption{Reserved failure cases: faint-sheet loss, residual blobs,
  cross-wrap risk, and examples whose morphology changes near the selected
  operating threshold.}
  \Description{Reserved examples of model and postprocess failure cases.}
  \label{fig:limitations}
\end{figure}

Finally, exact replay requires large, separately licensed artifacts.
The training code, recipes, hashes, and export scaffolding are public-repository
material, but the approximately 31 GB validation corpus, teacher atlases, M7
weights, and student checkpoints require dedicated artifact hosting and resolved
upstream terms.

\section{Conclusion}
\label{sec:conclusion}

We presented a cross-resolution supervision method for improving the released
M7 papyrus-surface segmenter without changing its deployed architecture.
A fine teacher exposes local inconsistencies through soft occupancy, medial
evidence, and de-blob separation; M7 preservation and a global trust region keep
the answer close to the incumbent function.
The same design principle extends downstream: the 2D fitter may propose a
connection, but radial, merger, cross-sheet, crossing, and persistence gates
force it to answer a sequence of geometric objections before adding pixels.

All measured intervals improve both calibration scrolls. The best
overlap diagnostic occurs at \bestSamplesText{}
samples, but locked and blinded morphology select the raw
\selectedSamplesText-sample checkpoint at $T=\selectedThreshold$. That choice
passes the anti-blob gate with no interior or thickness regressions while
retaining the longer run's longitudinal growth. The remaining release work is
to package the selected weights under resolved upstream terms and replace the
reserved panels with registered figures; the model decision itself is no
longer provisional.

\clearpage
\onecolumn
\section{Qualitative Gallery}
\label{sec:gallery}

The following pages collect the clearest qualitative results from the selected
raw checkpoint. They are presentation views, not additional samples or a hidden
selection stage: every crop appears in the committed locked-16 or fixed
PHerc1447 review corpus, and every student panel uses 8,192 samples at
$t=0.45$ without blending.

\begin{center}
  \figureasset{gallery_locked_growth}{0.86\textwidth}{7.4in}
  \captionsetup{hypcap=false}
  \captionof{figure}{Locked PHerc0139 growth gallery. Every registered crop is
  shown as CT, released M7, projected fine-teacher target, and raw M7-XR. These
  examples emphasize learned longitudinal regularity while keeping the released
  state of the art visible.}
  \label{fig:gallery-locked}
\end{center}

\clearpage
\begin{center}
  \figureasset{gallery_blind_antiblob}{0.86\textwidth}{7.4in}
  \captionsetup{hypcap=false}
  \captionof{figure}{Blind PHerc1447 anti-blob gallery. The released M7 foreground is cyan
  and the selected raw student's foreground is magenta, each on the identical
  CT slice. No fine teacher exists for these views.}
  \label{fig:gallery-blind}
\end{center}

\clearpage
\twocolumn


\bibliographystyle{ACM-Reference-Format}
\bibliography{submissions/2026-09/src/references}

\end{document}
